\question[4]
In Java werden ganze Zahlen mit einer festen Anzahl von Bits im Zweierkomplement gespeichert.
\begin{parts}
\part Gib den Wertebereich für 8 Bit (\texttt{byte}) und für 16 Bit (\texttt{short}) an.
\part Gegeben ist der Java-Code
\begin{lstlisting}[language=Java]
byte b = 127;
b++;
\end{lstlisting}
Danach hat \texttt{b} den Wert $-128$. Erläutere, wie das zustande kommt.
\part In Python ergibt \texttt{127 + 1} immer \texttt{128}. Erläutere, warum.
\end{parts}

\begin{solutionbox}{6cm}
\begin{parts}
\part 8 Bit: $-2^7 = -128$ bis $2^7 - 1 = 127$. 16 Bit: $-2^{15} = -32\,768$ bis $2^{15} - 1 = 32\,767$.
\part 127 ist der größte Wert, binär \texttt{01111111}. Plus 1 ergibt \texttt{10000000}. Im Zweierkomplement ist das
höchste Bit das Vorzeichen, \texttt{10000000} bedeutet $-128$. Der Wertebereich wird überschritten, es gibt einen
Überlauf ohne Fehlermeldung.
\part In Python hat \texttt{int} keine feste Anzahl von Bits. Eine ganze Zahl kann beliebig groß werden, solange der
Arbeitsspeicher reicht. Deshalb kann kein Überlauf passieren.
\end{parts}
\end{solutionbox}
